% Shared preamble of the task figures of Lecture 9 (based on Lecture2/task-figures/common.tex).
% The colours and the mapping are those of ../main.tex: Alice indigo, Mallory pink, public keys violet,
% private keys amber, Bob teal, the CA and its signature brown.
\usepackage{lmodern}
\usepackage[T1]{fontenc}
\renewcommand{\familydefault}{\sfdefault}
\usepackage{amsmath,amssymb}
\usepackage{xcolor}
\usepackage{fontawesome5}
\usetikzlibrary{arrows.meta, calc, positioning, fit, decorations.pathreplacing, decorations.pathmorphing, backgrounds}

% ---------- template ----------
\definecolor{DeepBlue}{HTML}{283593}
\definecolor{AccentOrange}{HTML}{EB5A14}
\definecolor{LightGray}{HTML}{F5F6FA}
% ---------- actors and categories (S1 to S6) ----------
\definecolor{sIndigo}{HTML}{5C6BC0}   % Alice
\definecolor{sPink}{HTML}{EC407A}     % Mallory
\definecolor{sViolet}{HTML}{7E57C2}   % public keys
\definecolor{sAmber}{HTML}{FFB300}    % private keys
\definecolor{sTeal}{HTML}{26A69A}     % Bob
\definecolor{sBrown}{HTML}{8D6E63}    % the CA and its signature
\definecolor{sPinkText}{HTML}{C2185B}
\definecolor{sAmberText}{HTML}{A06A00}
\definecolor{sTealText}{HTML}{00796B}
% ---------- actions and results ----------
\definecolor{aExec}{HTML}{43A047}
\definecolor{aRead}{HTML}{1E88E5}
\definecolor{aWrite}{HTML}{E53935}
% ---------- neutral ----------
\definecolor{nInk}{HTML}{37474F}
\definecolor{nLine}{HTML}{90A4AE}
\definecolor{nLight}{HTML}{CFD8DC}
\definecolor{nFill}{HTML}{ECEFF1}
\definecolor{nIdle}{HTML}{E0E0E0}

% ---------- TikZ styles (as in main.tex; byte cells a little larger for the web) ----------
\tikzset{
  hw/.style={draw=nLine, thick, fill=nFill, text=nInk, rounded corners=3pt, align=center, inner sep=4pt},
  hwmain/.style={draw=nInk, thick, fill=nInk, text=white, rounded corners=3pt, align=center, inner sep=4pt},
  cell/.style={draw=nLine, thick, fill=white, text=nInk, font=\ttfamily\small},
  cellfree/.style={cell, dashed, draw=nLight},
  container/.style={draw=nLine, thick, dashed, rounded corners=5pt},
  actor/.style={fill=#1, draw=#1, text=white, rounded corners=3pt, align=center, font=\small\bfseries, inner sep=5pt},
  actorlight/.style={fill=#1!15, draw=#1, thick, text=nInk, rounded corners=3pt, align=center, inner sep=5pt},
  data/.style={-{Stealth}, very thick, nInk},
  flow/.style={-{Stealth}, very thick, DeepBlue},
  annot/.style={draw=AccentOrange, line width=1.6pt, rounded corners=3pt},
  byte/.style={draw=nLine, thick, fill=white, text=nInk, font=\ttfamily\small,
    minimum width=0.8cm, minimum height=0.62cm, inner sep=1pt, rounded corners=1.5pt},
  ptb/.style={byte, fill=white, draw=nInk},   % public keys: neutral here, violet is the ciphertext
  ksb/.style={byte, fill=sAmber!25, draw=sAmber},
  rowlabel/.style={font=\small\bfseries, text=nInk, anchor=east},
  note/.style={font=\small, text=nInk},
  bigxor/.style={circle, draw=nInk, very thick, fill=white, inner sep=0pt, minimum size=0.8cm, font=\Large},
  radio/.style={-{Stealth}, very thick, nInk, decorate,
    decoration={snake, amplitude=1.2pt, segment length=7pt, post length=4pt}},
  % a white card behind the whole figure, so that it stays readable in a dark theme
  card/.style={show background rectangle, background rectangle/.style={fill=white, draw=none}, inner frame sep=8pt},
}
% a row of bytes: \byterow{x of the first cell}{y}{style}{list of values}{step between cells}
\newcommand{\byterow}[5]{%
  \expandafter\def\expandafter\brlist\expandafter{#4}%
  \foreach \v [count=\k from 0] in \brlist {
    \node[#3] at ({#1+#5*\k},#2) {\v};
  }}
% the generator / cipher box: dark, with cogs. \cogbox{name}{position}{label}{width}
\newcommand{\cogbox}[4]{%
  \node[hwmain, minimum width=#4, minimum height=1.3cm] (#1) at (#2) {};
  \node[text=white, font=\Large] at ([yshift=0.22cm]#1.center) {\faCogs};
  \node[text=white, font=\small\bfseries] at ([yshift=-0.32cm]#1.center) {#3};}
\newcommand{\ok}{\textcolor{aExec}{\faCheck}}
\newcommand{\nok}{\textcolor{aWrite}{\faTimes}}
\newcommand{\badge}[1]{\tikz[baseline=(b.base)]\node[circle,fill=AccentOrange,text=white,
  inner sep=1.5pt,font=\bfseries\footnotesize] (b) {#1};}


% ---------- Lecture 9 ----------
\tikzset{
  person/.style={circle, fill=#1, draw=#1, text=white, minimum size=1.3cm, font=\LARGE, inner sep=0pt},
  pubval/.style={draw=sViolet, thick, fill=sViolet!12, text=nInk, rounded corners=3pt, align=center, inner sep=4pt},
  privval/.style={draw=sAmber, thick, fill=sAmber!25, text=nInk, rounded corners=3pt, align=center, inner sep=4pt},
  caval/.style={draw=sBrown, thick, fill=sBrown!15, text=nInk, rounded corners=3pt, align=center, inner sep=4pt},
  cardbox/.style={draw=nLine, thick, fill=white, rounded corners=4pt},
  seal/.style={circle, fill=sBrown, draw=sBrown!70!black, line width=1pt, text=white, inner sep=0pt},
  todo/.style={draw=AccentOrange, line width=1.4pt, dashed, rounded corners=3pt, fill=AccentOrange!8, text=nInk, align=left, inner sep=5pt, font=\small},
  term/.style={draw=nLine, thick, fill=nFill, rounded corners=3pt, inner sep=6pt, font=\ttfamily\small, text=nInk, align=left},
}
\newcommand{\personnode}[4]{%
  \node[person=#3] (#1) at (#2) {\faUser};
  \node[font=\bfseries, text=#3, below=0.08cm of #1] (#1lbl) {#4};}
\newcommand{\mallorynode}[2]{%
  \node[person=sPink] (#1) at (#2) {\faUserSecret};
  \node[font=\bfseries, text=sPinkText, below=0.08cm of #1] (#1lbl) {Mallory};}
\newcommand{\canode}[3]{%
  \node[circle, fill=sBrown, text=white, minimum size=1.3cm, font=\LARGE, inner sep=0pt] (#1) at (#2) {\faLandmark};
  \node[font=\bfseries, text=sBrown, below=0.08cm of #1] (#1lbl) {#3};}
\newcommand{\five}{-{}-{}-{}-{}-}
